\section{Provider Taxonomy and Problem Formulation}
\label{sec:providers}

We model an inference marketplace as a set of services $S$ that expose the \emph{same} underlying LLMs under different pricing and capacity contracts.
Requests arrive online as a sequence $\{r_i\}$; request $r_i$ has input length $n_i^{\text{in}}$ and an \emph{a priori unknown} output length $n_i^{\text{out}}$.
Our routing decision chooses exactly one service for each request, trading off monetary cost against latency constraints.
\begin{table}[t]
\centering
\begin{tabular}{@{}lcl@{}}
\toprule
\textbf{Provider Type} & \textbf{Symbol} & \textbf{Examples} \\
\midrule
Pay-per-token API & $S_A$ & Together \\
Quota subscription & $S_Q$ & Chutes.ai \\
Concurrency subscription & $S_C$ & Featherless.ai \\
\bottomrule
\end{tabular}
\caption{Provider taxonomy for open-weight model inference (Llama, Mistral, DeepSeek-R1). Formal constraints are defined in Sections~\ref{sec:sa}--\ref{sec:sc}.}
\label{tab:provider_taxonomy}
\end{table}

\subsection{Pay-per-token APIs ($S_A$)}
\label{sec:sa}
A pay-per-token API charges a posted price per input and output token.
Routing request $r_i$ to $S_A$ incurs
\[
c_i \;=\; p_{\text{in}} \cdot n_i^{\text{in}} \;+\; p_{\text{out}} \cdot n_i^{\text{out}}.
\]
We use this service as the \emph{elastic overflow tier}: it is always feasible (no hard capacity constraint) but can be expensive at the margin.
In practice, multiple providers expose open-weight models under this pricing structure with different posted rates.\cite{together_pricing,fireworks_pricing,groq_pricing,deepseek_api_pricing}

\subsection{Quota subscriptions ($S_Q$)}
\label{sec:sq}
A quota subscription provides ``zero marginal cost'' inference up to a hard request budget within a fixed time window.
Formally, for each window, we require
\[
\sum_{i \in \text{window}} x_i^{Q} \;\le\; Q,
\]
where $x_i^{Q}\in\{0,1\}$ indicates whether request $i$ is routed to $S_Q$.
This model captures subscription tiers that impose request limits within fixed time windows (e.g., Chutes.ai with daily request quotas).\cite{chutes_pricing}
The key algorithmic challenge is \emph{opportunity cost}: allocating a scarce quota slot to a short/low-value request may force a later long/high-value request to the metered API.

\subsection{Concurrency subscriptions ($S_C$)}
\label{sec:sc}
A concurrency subscription exposes a reserved serving capacity that can process up to $C$ requests in parallel.
If request $i$ starts at time $s_i$ and occupies the server for processing time $p_i$, then the in-flight constraint is
\[
\sum_i \mathbb{I}\!\left(t \in [s_i,\, s_i+p_i]\right) \;\le\; C \quad \forall t.
\]
This captures dedicated inference endpoints with reserved capacity (e.g., Featherless.ai, Arli AI),\cite{featherless-plans,arli-pricing}
as well as self-hosted deployments in which the number of replicas bounds the number of concurrent requests.\cite{vllm_pagedattention}
Compared to $S_Q$, $S_C$ introduces scheduling structure because the resource is blocked for the request duration.

\subsection{Subscription amortization}
\label{sec:amortization}
In our experiments, we treat subscription fees as sunk costs (already paid) and optimize \emph{additional} spend on $S_A$ given fixed subscription capacities.
This matches common operational settings in which an organization has committed to subscriptions (or provisioned a cluster) and aims to maximize utilization while minimizing overflow to token-metered APIs.

\subsection{A unified abstraction beyond public APIs}
While our motivating examples come from the API marketplace, the same resource types arise in local GPU deployments.
\textbf{Concurrency as dedicated capacity:} $S_C$ maps directly to a GPU cluster with a fixed number of serving replicas, where $C$ is the maximum parallel requests.\cite{vllm_pagedattention}
\textbf{Quota as gutter pool:} $S_Q$ maps to the \textit{gutter pool}---spare GPU nodes maintained for high availability.\cite{sre_introduction_goog}
Since GPU nodes fail frequently due to hardware errors or network partitions,\cite{kokolis_reliability_2024} operators provision hot-standby capacity; routing inference to these idle nodes extracts value from sunk cost, subject to availability constraints that map to our quota model.
This unified view connects API routing to internal GPU cluster scheduling.

\subsection{Optimization objective}
\label{sec:objective}
Given a stream of requests $\mathcal{R} = \{r_1, \dots, r_T\}$, our goal is to minimize the total API cost:
\[
\min \sum_{i \in \mathcal{R}} x_i^A \cdot c_i
\]
subject to the quota constraint $\sum_{i \in \text{window}} x_i^Q \le Q$ and concurrency constraint $\sum_i \mathbb{I}(t \in [s_i, s_i+p_i]) \le C$ for all $t$, where $x_i^A, x_i^Q, x_i^C \in \{0, 1\}$ indicate routing decisions and exactly one must be 1 for each request.

\subsection{Key assumptions}
\label{sec:assumptions}
We make two simplifying assumptions.
First, \textbf{homogeneous quality}: all providers serve the same model, so routing decisions affect only cost, not output quality. This holds for open-weight models available across multiple providers.
Second, the \textbf{output token count} $n_i^{\text{out}}$ is unknown at routing time---the key source of uncertainty that motivates our learning-augmented approach.
